% RQ3: pooling over a track (task S3, Person A).
% Every number below comes from results/rq3_pooling.json through
% figures/rq3_numbers.tex, which make_figures.py writes. A number that does not
% exist yet prints as ?? in bold. Lines marked TODO need the results: write
% them once scripts/research/pooling_eval.py has run on the labelled clips.

\providecommand{\rqthree}[2]{%
  \ifcsname rq3/#1/#2\endcsname\csname rq3/#1/#2\endcsname\else\textbf{??}\fi}
\InputIfFileExists{figures/rq3_numbers.tex}{}{}

\section{RQ3: Which Frames Should Count?}
\label{sec:rq3}
A person in a store stays in view for many frames, and the frames differ. Some
crops are sharp and whole. Others are blurred by motion, cut off at the edge of
the frame or half hidden behind another shopper. The majority vote of
Eq.~\eqref{eq:majority} gives every crop the same say. RQ3 asks whether a rule
that trusts the good crops more names more people correctly, and whether hand
measured image quality does as well as a weighting learned from video labels.
Our hypothesis (H3) is that quality weighting beats the majority vote, with the
largest gain on people who are truncated or occluded.

\subsection{Pooling Rules}
Apart from the majority vote, every rule takes a weighted mean of the frames'
probability vectors and picks its largest entry:
\begin{equation}
\bar{p}(c) = \frac{\sum_i w_i\, p_i(c)}{\sum_i w_i}, \qquad
\hat{y} = \arg\max_{c} \bar{p}(c).
\label{eq:weighted}
\end{equation}
The rules differ only in the weights $w_i$.
\begin{itemize}
  \item \emph{Single frame.} One frame chosen at random. This is what a system
        without pooling would report.
  \item \emph{Majority vote.} Eq.~\eqref{eq:majority}, the product's rule.
  \item \emph{Mean probability.} $w_i = 1$.
  \item \emph{Confidence weighted.} $w_i = \max_c p_i(c)$, the classifier's own
        confidence in the frame.
  \item \emph{Quality weighted.} $w_i = \Pr(\mathrm{correct}_i \mid q_i)$, the
        chance that frame $i$ is classified correctly given its cues, from the
        quality model below.
\end{itemize}

\subsection{Quality Cues and the Quality Model}
The cue vector $q_i$ has seven entries, all measured when the crop is exported.
\begin{itemize}
  \item \emph{Sharpness.} The share of the crop's spectral energy at or above
        0.25 cycles per pixel, computed on a $128\times128$ grey version. Blur
        and low resolution both remove this band. The variance of the
        Laplacian~\cite{pech2000laplacian} is kept as a second sharpness cue.
  \item \emph{Motion.} Dense Farneb\"ack optical flow~\cite{farneback2003flow}
        between the crop's frame and the previous one: the mean flow magnitude
        inside the box minus the median magnitude outside every box, in pixels
        per frame. Subtracting the background removes camera shake, so what
        remains is the person's own movement, which is what blurs the crop.
  \item \emph{Truncation.} 1 if the box touches the frame border.
  \item \emph{Occlusion.} The box's largest IoU with another person's box in the
        same frame.
  \item \emph{Size and detector confidence.} The box's share of the frame area
        and the detector's score.
\end{itemize}
The quality model is a logistic regression with an L2 penalty from these cues
to whether the classifier got the crop right. The Laplacian variance and the
size enter as logarithms and the motion as $\log(1+m)$, because each spans
several orders of magnitude. Features are standardised, a cue that could not be
measured takes the training mean, and the model is fitted by Newton's method.
It is small and readable, and its coefficients show which cue matters
(Section~\ref{sec:rq3-results}).

\subsection{Learned Attention Baseline}
Video attribute recognition usually learns how to weigh frames from video
labels~\cite{chen2019temporal}. A reviewer will rightly ask how measured quality
compares with that, so we train the same kind of head on top of the frozen
per-frame features of our two classifiers: the 768-dimensional ViT feature that
the gender head reads, and the histograms that the age SVMs score. With $f_i$
the standardised feature of frame $i$,
\begin{gather}
z_i = W f_i + b, \qquad
\alpha_i = \frac{\exp\left(v^{\top}\tanh z_i\right)}
                {\sum_j \exp\left(v^{\top}\tanh z_j\right)}, \\
\hat{y} = \arg\max_c \Big[\, U \sum\nolimits_i \alpha_i z_i + u \,\Big]_c,
\end{gather}
with $z_i \in \mathbb{R}^{32}$. The head is trained with class-balanced
cross-entropy and Adam for 300 epochs with a weight decay of $10^{-3}$. It
needs a label for every training person. Quality weighting needs only crops
labelled right or wrong, and that difference in what each method asks for is
part of the comparison. MARS~\cite{zheng2016mars} with its tracklet
attributes~\cite{chen2019temporal} is the natural place to train the head. When
MARS is not available we train it with the same cross-validation folds as the
quality model.

\subsection{Protocol}
Two annotators labelled every person in the retail test clips from the crops
alone, without seeing the model's answers, and resolved their disagreements
together. Agreement before resolution is reported as Cohen's
$\kappa$~\cite{cohen1960kappa}. People either annotator marked unsure are left
out, and their number is reported.

The per-crop probabilities are the product's own, stored when the crops were
exported, so every rule pools exactly what the deployed system saw. The quality
model and the attention head are fitted with 5-fold cross-validation grouped by
person: no person's crops ever help weigh that person's own frames. The unit of
evaluation is the person. We report per-person accuracy with a 95\% bootstrap
interval over people~\cite{efron1993bootstrap} (2,000 resamples), and McNemar's
test~\cite{mcnemar1947} of each rule against the majority vote on the same
people, exact when fewer than 25 people disagree. To see how much pooling is
worth, we also draw $k = 1, \dots, 12$ of a person's crops at random (20 draws
per person) and pool only those. Last, we split people into those who are
truncated, or occluded, in at least half their crops, and the rest.

\begin{table}[t]
\caption{Per-person accuracy of each pooling rule with its 95\% interval over
people, and McNemar's $p$ against the majority vote. Gender: \rqthree{gender}{n}
people, age: \rqthree{age}{n} people.}
\label{tab:rq3}
\centering
\begin{tabular}{@{}lcccc@{}}
\toprule
Rule & Gender & $p$ & Age & $p$ \\
\midrule
Single frame        & \rqthree{gender}{single}     & \rqthree{gender}{p-single}     & \rqthree{age}{single}     & \rqthree{age}{p-single} \\
Majority vote       & \rqthree{gender}{majority}   & --                             & \rqthree{age}{majority}   & -- \\
Mean probability    & \rqthree{gender}{mean}       & \rqthree{gender}{p-mean}       & \rqthree{age}{mean}       & \rqthree{age}{p-mean} \\
Confidence weighted & \rqthree{gender}{confidence} & \rqthree{gender}{p-confidence} & \rqthree{age}{confidence} & \rqthree{age}{p-confidence} \\
Quality weighted    & \rqthree{gender}{quality}    & \rqthree{gender}{p-quality}    & \rqthree{age}{quality}    & \rqthree{age}{p-quality} \\
Learned attention   & \rqthree{gender}{attention}  & \rqthree{gender}{p-attention}  & \rqthree{age}{attention}  & \rqthree{age}{p-attention} \\
\bottomrule
\end{tabular}
\end{table}

\begin{figure*}[t]
\centering
\IfFileExists{figures/rq3_pooling.pdf}{%
  \includegraphics[width=\textwidth]{figures/rq3_pooling.pdf}}{%
  \fbox{\parbox{0.9\textwidth}{\centering\vspace{6mm}This figure is drawn by
  make\_figures.py once results/rq3\_pooling.json exists.\vspace{6mm}}}}
\caption{RQ3 on the retail test clips. (a) Per-person accuracy of each pooling
rule with its 95\% interval over people. (b) Accuracy when only $k$ of a
person's crops are pooled.}
\label{fig:rq3}
\end{figure*}

\subsection{Results}
\label{sec:rq3-results}
Table~\ref{tab:rq3} and Fig.~\ref{fig:rq3} give the results. With no pooling,
a single frame gave a per-person gender accuracy of
\rqthree{gender}{single-value}. The majority vote reached \rqthree{gender}{majority-value} and quality weighting
\rqthree{gender}{quality-value} (McNemar $p = \rqthree{gender}{p-quality}$, with
\rqthree{gender}{b-quality} people right only under quality weighting and
\rqthree{gender}{c-quality} right only under the vote). For age the same three
numbers were \rqthree{age}{single-value}, \rqthree{age}{majority-value} and
\rqthree{age}{quality-value}.
% TODO after results: one sentence on whether H3 holds for gender and for age.

On the \rqthree{gender}{truncated-n} people who were truncated in most of their
crops, the vote reached \rqthree{gender}{truncated-majority} and quality
weighting \rqthree{gender}{truncated-quality}. On the \rqthree{gender}{occluded-n}
occluded people the two reached \rqthree{gender}{occluded-majority} and
\rqthree{gender}{occluded-quality}.
% TODO after results: is the gain larger on these subsets, as H3 predicts?

The learned attention head reached \rqthree{gender}{attention-value} on gender
and \rqthree{age}{attention-value} on age, trained with
\rqthree{gender}{attention-fitted}. Its frame weights and the quality weights had a mean
Spearman correlation of \rqthree{gender}{rho} for gender.
% TODO after results: if rho is high, the head learned roughly what the cues
% measure; if low, it attends to something else. Say which, and what follows
% for a deployment that has no video labels.

In the gender quality model fitted on all people, the standardised coefficients
were \rqthree{gender}{coef-hf_ratio} for sharpness,
\rqthree{gender}{coef-log1p_motion} for motion, \rqthree{gender}{coef-truncated}
for truncation and \rqthree{gender}{coef-occlusion} for occlusion. A positive
coefficient means crops with more of that cue are classified correctly more
often.
% TODO after results: name the cue that matters most, and say if any sign is
% the opposite of what one would expect.

Fig.~\ref{fig:rq3}b shows how accuracy grows with the number of frames pooled.
% TODO after results: where the curve flattens. If it flattens well before 12,
% the vote cap could be lowered and the classifier run less often.
